% Gaussian-posterior part of the exponential sample lower-bound theorem.
% Requires theorem, lemma, proposition, corollary environments and amsmath,
% amssymb, amsthm. The spectral/random-Gram lemma is supplied separately.

\section{The actual Gaussian sign posterior}
Let $\mathcal X=\{-1,1\}^d$, $M=2^d$, and
$f(x)=\prod_{j\in U}x_j, \qquad |U|=k$. All function norms and expectations on
$\mathcal X$ use its uniform probability measure. Let $G$ be a
nondegenerate centered Gaussian vector indexed by $\mathcal X$, with
covariance matrix $K$. The training inputs $X_1,\ldots,X_n$ are independent
and uniform. We condition on
\[
 \mathcal A_S=\{f(X_i)G(X_i)>0:1\le i\le n\}.
\]
The prediction at $x$ is its actual conditional correct-label probability
$q_S(x)=\mathbb P(f(x)G(x)>0\mid\mathcal A_S)$, and we write
\[
 F_\theta(S)=\mathbb P_x\{q_S(x)\le\theta\},\qquad
 \theta=\frac{13}{20},\qquad
 c_\theta=\sqrt{2\pi}\left(\frac1e-\frac7{20}\right)>0.
\]
The threshold $\theta<2/3$ supplies a strict buffer for passage from
limiting Gaussian posteriors to finite-width posterior probabilities.

Suppose the Walsh characters are eigenfunctions of the normalized
covariance operator, with degree-$j$ eigenvalue $\lambda_j$. Write
$V_r=\operatorname{span}\{\chi_U:|U|\le r\}$ and $D_r=\dim V_r$.
For $r<d$, let $Q$ be the covariance contributed by the complementary
Fourier modes, and put
\[
 \tau=Q(x,x)=\sum_{j>r}\binom dj\lambda_j,
 \qquad \ell=\max_{j>r}\lambda_j.
\]
The matrix eigenvalues of $Q$ are $M\lambda_j$; this normalization is
used throughout. Define $\|h\|_K^2=h^TK^{-1}h$, and write $u_+=\max(u,0)$ and $u_-=\max(-u,0)$.

\begin{lemma}[Mean norm and a confidence margin]\label{lem:schedule-mean}
Let $m(x)=\mathbb E[G(x)\mid\mathcal A_S]$, and let $s\le n$ be the
number of distinct sampled inputs. Then
\[
 \|m\|_K^2\le s,\qquad f(X_i)m(X_i)>0.
\]
Let $\sigma_x^2=\operatorname{Var}(G(x)\mid G_S)$, where $S$ here lists
the distinct inputs. Whenever $\sigma_x>0$ and $q_S(x)>\theta$,
\[
 f(x)m(x)>c_\theta\sigma_x.
\]
\end{lemma}
\begin{proof}
Whiten the Gaussian training vector as $G_S=K_{SS}^{1/2}\xi$.
The sign constraints form a homogeneous cone in $\xi$. Conditioning
leaves the squared Gaussian radius with mean $s$, so Jensen's inequality
and Gaussian regression give
\[
 m_S^TK_{SS}^{-1}m_S\le s,\qquad
 m=K_{\mathcal X,S}K_{SS}^{-1}m_S.
\]
The RKHS norm of this regression vector is the preceding quadratic
form. Its observed signed coordinates are strictly positive.

For the confidence claim, the scalar conditional variable
$Y=f(x)G(x)$ has a log-concave density: it is a marginal of a Gaussian
restricted to a convex set, by the marginal log-concavity theorem \cite{prekopa1973}. If $\mu=\mathbb E Y$ and $F_Y$ is its
cdf, log-concavity and Jensen's inequality imply
\[
 \log F_Y(\mu)\ge\mathbb E\log F_Y(Y)=-1.
\]
Consequently $\mathbb P(Y>\mu)\le1-e^{-1}<\theta$.
Conditioning first on the full training vector shows that the density
of $Y$ is a mixture of Gaussian densities with common variance
$\sigma_x^2$, and is therefore at most
$(\sqrt{2\pi}\sigma_x)^{-1}$. If $q_S(x)>\theta$, then $\mu>0$ and
\[
 \theta<\mathbb P(Y>0)
 \le 1-e^{-1}+\frac{\mu}{\sqrt{2\pi}\sigma_x}.
\]
This gives the stated margin.
\end{proof}

\begin{lemma}[Most residual variances remain large]\label{lem:schedule-var}
For every training sample,
\[
 \mathbb P_x\{\sigma_x^2<\tau/2\}\le \frac{2n\ell}{\tau}.
\]
Also, if $p$ is the degree-$r$ Fourier truncation of $m$ and $t=m-p$,
then $\|t\|_2^2\le n\ell$.
\end{lemma}
\begin{proof}
Write $G=G_{\le r}+G_{>r}$ as a sum of independent fields. Conditioning
additionally on the full low-degree field can only decrease conditional
Gaussian variance. The remaining covariance is
\[
 Q-Q_{\mathcal X,S}Q_{SS}^{\dagger}Q_{S,\mathcal X}.
\]
Its subtracted term is positive semidefinite, bounded above by $Q$,
and has rank at most $s\le n$. Its normalized trace is at most $n\ell$.
Markov's inequality proves the first assertion; it includes observed
points automatically. Fourier orthogonality and the mean norm bound
prove $\|t\|_2^2\le\ell\|m\|_K^2\le n\ell$.
\end{proof}

\section{Uniform sampling detects the negative part}
\begin{lemma}[Uniform ReLU sampling on the cube]\label{lem:schedule-relu}
For every fixed $\delta,\eta\in(0,1)$, there is a universal constant
$C$ such that, if
\[
 n\ge C\eta^{-2}D_r\log\!\left(\frac{D_r}{\delta\eta}\right),
\]
then with probability at least $1-\delta$, simultaneously for all
$p\in V_r$,
\[
 \left|\frac1n\sum_{i=1}^n(-f(X_i)p(X_i))_+
             -\mathbb E_x(-f(x)p(x))_+\right|
 \le\eta\mathbb E_x|p(x)|.
\]
In particular, when $r<k$ and $\eta=1/4$, the empirical negative
part is at least $\mathbb E|p|/4$.
\end{lemma}
\begin{proof}
Let $A$ have rows $f(x)(\chi_U(x))_{|U|\le r}$, one for each cube
point. Then $A^TA=MI$ and every row has squared norm $D_r$.
The $\ell_1$ Lewis weights of $A$ are exactly $D_r/M$: substitution
in their defining equation
\[
 w_x^2=A_x(A^T\operatorname{diag}(w^{-1})A)^{-1}A_x^T
\]
verifies this constant solution, and Lewis weights are unique.
The additive ReLU coreset theorem of Mai, Musco, and Rao
\cite[Theorem~5]{mai2021coresets} therefore samples precisely the uniform
distribution on the cube and gives the stated estimate. Its sample bound
is $O(\eta^{-2}D_r\log(D_r/(\delta\eta)))$; no additional logarithmic
factor is required. Finally $\mathbb E fp=0$, because $r<k$, so
$\mathbb E(-fp)_+=\mathbb E|p|/2$.
\end{proof}

